% =====================================================================
% VERIFICA DI MATEMATICA - 2BS
% 60 minuti, 30 punti.
% Punteggi: 4 + 8 + 10 + 3 + 5 = 30.
% Compilare almeno due volte con pdfLaTeX.
% Per mostrare le soluzioni: decommentare \printanswers.
% =====================================================================

\documentclass[11pt,a4paper,addpoints]{exam}

% ---------- Lingua e codifica ----------
\usepackage[utf8]{inputenc}
\usepackage[T1]{fontenc}
\usepackage[italian]{babel}
\usepackage{lmodern}
\usepackage{microtype}

% ---------- Matematica e impaginazione ----------
\usepackage{amsmath,amssymb}
\usepackage{xcolor}
\usepackage{array,tabularx}
\usepackage[
  a4paper,
  top=2cm,
  bottom=2.2cm,
  left=2cm,
  right=2.2cm,
  headheight=30pt,
  headsep=12pt
]{geometry}

% ---------- Soluzioni ----------
%\printanswers
\ifdefined\CORRETTORE
  \printanswers
\fi
\cancelspace
\unframedsolutions
\renewcommand{\solutiontitle}{%
  \noindent\textbf{Soluzione.}\enspace
}
\SolutionEmphasis{\color{blue!60!black}}

% ---------- Dati della prova ----------
\newcommand{\materia}{Matematica}
\newcommand{\classe}{2BS}
\newcommand{\indirizzo}{Liceo Scientifico}
\newcommand{\tipoprova}{Verifica di matematica}
\newcommand{\argomento}{Frazioni algebriche ed equazioni frazionarie}
\newcommand{\durata}{60 minuti}
\newcommand{\annoscolastico}{2026/2027}
\newcommand{\docente}{Prof. Rubens Longhi}
\newcommand{\fila}{A}

\newcommand{\indicazioni}{%
  Svolgi sul foglio a protocollo e mostra i passaggi.
  Indica sempre le condizioni di esistenza prima di semplificare.
  Per le equazioni scrivi l'insieme delle soluzioni.
  Lavora in $\mathbb{R}$, senza calcolatrice.
}

% ---------- Macro matematiche ----------
\newcommand{\R}{\mathbb{R}}

% ---------- Intestazioni e piè di pagina ----------
\pagestyle{headandfoot}
\firstpageheader
  {%
    \scriptsize\textsc{IIS David Maria Turoldo}\\[-1pt]
    \scriptsize Via Ronco 11 -- 24019 Zogno (BG)
  }
  {}
  {%
    \scriptsize A.S.\ \annoscolastico\\[-1pt]
    \scriptsize \docente
  }
\firstpageheadrule

\runningheader
  {\small\materia\ -- \classe}
  {\small\tipoprova}
  {\small Fila \fila}
\runningheadrule

\newcommand{\paginefila}{\pageref{fine@fila@\fila}}

\firstpagefooter{}{\small Pagina \thepage\ di \paginefila}{}
\runningfooter{}{\small Pagina \thepage\ di \paginefila}{}

% ---------- Obiettivi e formato degli esercizi ----------
\newcommand{\obiettivocorrente}{}

\newcommand{\obiettivo}[1]{%
  \gdef\obiettivocorrente{#1}%
}

\newcommand{\stampaobiettivo}{%
  \ifx\obiettivocorrente\empty
  \else
    {\footnotesize\color{black!55}%
      \textit{Obiettivo: \obiettivocorrente}}%
  \fi
}

\qformat{%
  \textbf{\thequestion.}\enspace
  \stampaobiettivo
  \hfill
  {\small(\totalpoints\ \points)}%
  \gdef\obiettivocorrente{}%
}

\renewcommand{\questionlabel}{}
\renewcommand{\partlabel}{\thequestion\thepartno)}
\renewcommand{\subpartlabel}{\thesubpart.}

\pointpoints{punto}{punti}

\renewcommand{\partshook}{%
  \setlength{\itemsep}{4pt}%
  \setlength{\rightmargin}{0.5cm}%
}

% ---------- Numerazione indipendente per ciascuna fila ----------
\newcounter{offsetfila}
\newcounter{totalqfila}

\renewcommand{\thequestion}{%
  \number\numexpr
    \value{question}-\value{offsetfila}%
  \relax
}

\newenvironment{esercizi}
  {%
    \begin{questions}
    \setcounter{question}{\value{offsetfila}}%
  }
  {\end{questions}}

% ---------- Esercizi non spezzati nella versione studenti ----------
\newenvironment{esercizio}
  {%
    \par
    \ifprintanswers\else\begin{samepage}\fi
  }
  {%
    \ifprintanswers\else\end{samepage}\fi
    \par\vspace{3pt}
  }

% ---------- Tabella dei punteggi ----------
\hqword{\textbf{Esercizio}}
\hpword{\textbf{Punti massimi}}
\hsword{\textbf{Punti ottenuti}}
\htword{\textbf{Totale}}

% ---------- Inizio e fine di ciascuna fila ----------
\newcommand{\iniziofila}[1]{%
  \clearpage
  \renewcommand{\fila}{#1}%
  \setcounter{page}{1}%
  \setcounter{offsetfila}{\value{totalqfila}}%

  \begin{center}
    {\small\bfseries
      \tipoprova
      \ifprintanswers
        -- Soluzioni e criteri di correzione
      \fi
    }\\[3pt]
    {\Large \argomento}\\[2pt]
    {\small
      \indirizzo\ -- Classe \classe\
      -- Durata: \durata\
      -- \textbf{Fila #1}
    }
  \end{center}

  \vspace{4pt}
  \noindent
  \begin{tabularx}{\textwidth}{@{}l X l p{3cm}@{}}
    \textbf{Nome e Cognome:} & \hrulefill
    & \textbf{Data:} & \hrulefill\\[8pt]
  \end{tabularx}

  \noindent
  \fbox{%
    \parbox{\dimexpr\textwidth-2\fboxsep-2\fboxrule}{%
      \small
      \textbf{Indicazioni.} \indicazioni\space
      Punteggio totale: \pointsinrange{fila#1} punti.
    }%
  }

  \medskip
  \begingradingrange{fila#1}%
}

\newcommand{\finefila}{%
  \setcounter{totalqfila}{\value{question}}%
  \endgradingrange{fila\fila}%

  \ifprintanswers
  \else
    \vfill
    \begin{center}
      \small
      \partialgradetable{fila\fila}[h][questions]
    \end{center}

    \vspace{6pt}
    \noindent
    {\small
      \textbf{Voto:}\ \rule{3cm}{0.4pt}
      \hfill
      \textbf{Firma del docente:}\ \rule{5cm}{0.4pt}
    }%
  \fi

  \label{fine@fila@\fila}%
}

% =====================================================================
% DOCUMENTO
% =====================================================================
\begin{document}

\linespread{1.05}\selectfont
\ifprintanswers\small\fi

\iniziofila{A}
\begin{esercizi}

% ---------------- ESERCIZIO 1: 4 punti ----------------
\begin{esercizio}
\obiettivo{riconoscere proprietà e condizioni di validità}
\question
Stabilisci se ciascuna affermazione è vera o falsa.
Giustifica brevemente ogni risposta.

\begin{parts}
  \part[1]
  Le frazioni algebriche
  $\dfrac{x-4}{x+2}$ e $\dfrac{4-x}{-x-2}$
  sono equivalenti per tutti i valori per cui esistono.

  \begin{solution}
    \textbf{Vera.} Le C.E. sono $x\ne-2$ per entrambe.
    Cambiando segno sia al numeratore sia al denominatore
    il valore della frazione non cambia:
    \[
      \frac{4-x}{-x-2}
      =\frac{-(x-4)}{-(x+2)}
      =\frac{x-4}{x+2}.
    \]

    \textit{Criteri:} scelta 0,5 punti;
    giustificazione 0,5 punti.
  \end{solution}

  \part[1]
  Per dividere due frazioni algebriche è sufficiente
  imporre che i loro denominatori siano diversi da zero.

  \begin{solution}
    \textbf{Falsa.} Occorre anche che la frazione divisore
    sia diversa da zero: il suo numeratore deve essere non nullo.
    Per esempio, $\frac1x:\frac{x-1}{x+2}$ richiede
    $x\ne0$, $x\ne-2$ e anche $x\ne1$.

    \textit{Criteri:} scelta 0,5 punti;
    condizione aggiuntiva corretta 0,5 punti.
  \end{solution}

  \part[1]
  Nell'incognita $x$, l'equazione
  $\dfrac{3x}{a+2}=\dfrac{a-1}{a}$
  è frazionaria.

  \begin{solution}
    \textbf{Falsa.} È letterale e intera rispetto a $x$:
    l'incognita non compare nei denominatori.
    Il parametro deve soddisfare $a\ne-2$ e $a\ne0$.

    \textit{Criteri:} scelta 0,5 punti;
    spiegazione riferita all'incognita 0,5 punti.
  \end{solution}

  \part[1]
  Se, dopo aver imposto le C.E. ed eliminato i denominatori,
  un'equazione frazionaria si riduce a un'identità,
  le sue soluzioni sono tutti e soli i valori ammessi dalle C.E.

  \begin{solution}
    \textbf{Vera.} Nel dominio di esistenza il denominatore
    comune è non nullo, quindi la sua eliminazione è
    un passaggio equivalente. L'identità è verificata da
    ogni valore di tale dominio; i valori esclusi restano esclusi.

    \textit{Criteri:} scelta 0,5 punti;
    giustificazione 0,5 punti.
  \end{solution}
\end{parts}
\end{esercizio}

% ---------------- ESERCIZIO 2: 8 punti ----------------
\begin{esercizio}
\obiettivo{operare con frazioni algebriche e potenze negative}
\question
Determina le C.E. e semplifica le seguenti espressioni.

\begin{parts}
  \part[4]
  \[
    \left[
      \frac{x}{x^2-5x+6}
      +\frac{1}{x-2}
      -\frac{1}{x-3}
    \right]
    :
    \frac{2x-2}{x^2-4}.
  \]

  \begin{solution}
    Scomponiamo:
    \[
      x^2-5x+6=(x-2)(x-3),\qquad
      x^2-4=(x-2)(x+2).
    \]

    I denominatori devono essere non nulli e il divisore
    deve essere diverso da zero. Pertanto
    \[
      \boxed{\text{C.E.: }x\ne-2,\ 1,\ 2,\ 3}.
    \]

    La parentesi quadra diventa
    \[
      \frac{x+(x-3)-(x-2)}{(x-2)(x-3)}
      =\frac{x-1}{(x-2)(x-3)}.
    \]

    Quindi, mantenendo tutte le C.E. iniziali,
    \[
      \frac{x-1}{(x-2)(x-3)}
      \cdot\frac{(x-2)(x+2)}{2(x-1)}
      =\boxed{\frac{x+2}{2(x-3)}}.
    \]

    \textit{Criteri:} C.E. complete 1 punto;
    somma algebrica 1 punto;
    divisione e semplificazioni 1,5 punti;
    risultato 0,5 punti.
  \end{solution}

  \part[4]
  \[
    \left[
      \left(\frac{a^{-1}-3a^{-2}}{a^{-1}+3a^{-2}}\right)^2
      +2\frac{a^{-1}-3a^{-2}}{a^{-1}+3a^{-2}}
      +1
    \right]
    :
    \frac{2a^2}{a^3+6a^2+9a}.
  \]

  \begin{solution}
    Le potenze negative richiedono $a\ne0$. Inoltre
    \[
      a^{-1}+3a^{-2}=\frac{a+3}{a^2},\qquad
      a^3+6a^2+9a=a(a+3)^2.
    \]

    Il divisore è non nullo per $a\ne0$; si ottiene
    \[
      \boxed{\text{C.E.: }a\ne0,\ -3}.
    \]

    Nel dominio ammesso,
    \[
      \frac{a^{-1}-3a^{-2}}{a^{-1}+3a^{-2}}
      =\frac{(a-3)/a^2}{(a+3)/a^2}
      =\frac{a-3}{a+3}.
    \]

    Riconosciamo il quadrato di un binomio:
    \[
      \left(\frac{a-3}{a+3}\right)^2
      +2\frac{a-3}{a+3}+1
      =\left(\frac{a-3}{a+3}+1\right)^2
      =\frac{4a^2}{(a+3)^2}.
    \]

    Pertanto
    \[
      \frac{4a^2}{(a+3)^2}
      \cdot\frac{a(a+3)^2}{2a^2}
      =\boxed{2a}.
    \]

    Restano valide tutte le C.E. iniziali.

    \textit{Criteri:} C.E. 1 punto;
    gestione delle potenze negative 1 punto;
    parentesi quadra 1 punto;
    divisione e risultato 1 punto.
  \end{solution}
\end{parts}
\end{esercizio}

\ifprintanswers\else\newpage\fi

% ---------------- ESERCIZIO 3: 10 punti ----------------
\begin{esercizio}
\obiettivo{risolvere equazioni e controllare l'ammissibilità}
\question
Risolvi le seguenti equazioni, indicando le C.E.
e l'insieme delle soluzioni.

\begin{parts}
  \part[3]
  \[
    \frac{x^2+x-12}{x^2-9}=0.
  \]

  \begin{solution}
    \[
      x^2-9=(x-3)(x+3),\qquad
      \boxed{\text{C.E.: }x\ne-3,\ 3}.
    \]

    Una frazione è nulla se il numeratore è nullo
    e il denominatore è diverso da zero:
    \[
      x^2+x-12=(x+4)(x-3)=0.
    \]

    I candidati sono $x=-4$ e $x=3$.
    Il secondo è escluso dalle C.E., dunque
    \[
      \boxed{S=\{-4\}}.
    \]

    \textit{Criteri:} C.E. 1 punto;
    scomposizione e candidati 1 punto;
    selezione e insieme delle soluzioni 1 punto.
  \end{solution}

  \part[3]
  \[
    -\frac2x+\frac6{x+1}=1.
  \]

  \begin{solution}
    \[
      \boxed{\text{C.E.: }x\ne0,\ -1}.
    \]

    Moltiplicando entrambi i membri per $x(x+1)$:
    \[
      -2(x+1)+6x=x(x+1).
    \]

    Quindi
    \[
      x^2-3x+2=0
      \quad\Longleftrightarrow\quad
      (x-1)(x-2)=0.
    \]

    Entrambi i valori ottenuti sono ammessi:
    \[
      \boxed{S=\{1,2\}}.
    \]

    \textit{Criteri:} C.E. 0,5 punti;
    eliminazione dei denominatori e riduzione 1 punto;
    scomposizione 1 punto;
    insieme delle soluzioni 0,5 punti.
  \end{solution}

  \part[4]
  \[
    \left(
      \frac{2x+5}{x^2-9}+\frac1{3-x}
    \right)
    :
    \left(\frac2{9-x^2}\right)
    =-\frac52.
  \]

  \begin{solution}
    Tutti i denominatori devono essere non nulli.
    Il divisore, quando esiste, non è mai zero:
    \[
      \boxed{\text{C.E.: }x\ne-3,\ 3}.
    \]

    Poiché $3-x=-(x-3)$,
    \[
      \frac{2x+5}{(x-3)(x+3)}-\frac1{x-3}
      =\frac{2x+5-(x+3)}{(x-3)(x+3)}
      =\frac{x+2}{(x-3)(x+3)}.
    \]

    Inoltre $9-x^2=-(x-3)(x+3)$, quindi
    \[
      \frac{x+2}{(x-3)(x+3)}
      \cdot\frac{-(x-3)(x+3)}2
      =-\frac52
      \quad\Longleftrightarrow\quad
      -\frac{x+2}{2}=-\frac52.
    \]

    Si ottiene $x=3$, escluso dalle C.E.
    L'equazione è impossibile:
    \[
      \boxed{S=\varnothing}.
    \]

    \textit{Criteri:} C.E. 1 punto;
    somma e gestione dei segni 1 punto;
    divisione e riduzione 1 punto;
    esclusione del candidato e conclusione 1 punto.
  \end{solution}
\end{parts}
\end{esercizio}

% ---------------- ESERCIZIO 4: 3 punti ----------------
\begin{esercizio}
\obiettivo{ricavare una formula inversa}
\question[3]

Un veicolo percorre due tratti di uguale lunghezza
con velocità costanti positive $v_1$ e $v_2$.
La velocità media sull'intero percorso è
\[
  v_m=\frac{2v_1v_2}{v_1+v_2}.
\]

Ricava $v_2$ in funzione di $v_m$ e $v_1$.
Sapendo che $v_m>0$ e $v_1>0$, indica anche
quale condizione devono soddisfare affinché
il valore ricavato di $v_2$ sia positivo e finito.

\begin{solution}
  Poiché $v_1+v_2>0$, possiamo moltiplicare:
  \[
    v_m(v_1+v_2)=2v_1v_2
    \quad\Longleftrightarrow\quad
    v_mv_1=(2v_1-v_m)v_2.
  \]

  Se $2v_1-v_m\ne0$,
  \[
    \boxed{v_2=\frac{v_mv_1}{2v_1-v_m}}.
  \]

  Il numeratore è positivo. Per avere $v_2>0$ e finita
  occorre e basta che il denominatore sia positivo:
  \[
    \boxed{0<v_m<2v_1}.
  \]

  Per $v_m=2v_1$ non esiste una $v_2$ finita
  che soddisfi la relazione con $v_1>0$.

  \textit{Criteri:} passaggi e formula inversa 2 punti;
  condizione di validità fisica 1 punto.
\end{solution}
\end{esercizio}

% ---------------- ESERCIZIO 5: 5 punti ----------------
\begin{esercizio}
\obiettivo{interpretare e confrontare modelli frazionari}
\question

Le temperature di due componenti di una macchina,
in gradi Celsius, sono descritte dai modelli
\[
  T_1(t)=\frac{60(t+1)}{t+3},
  \qquad
  T_2(t)=\frac{30(2t+1)}{t+2},
\]
dove $t\ge0$ è il numero di minuti trascorsi dall'accensione.

\begin{parts}
  \part[1]
  Quali sono le temperature dei due componenti
  all'istante dell'accensione?

  \begin{solution}
    \[
      T_1(0)=\frac{60}{3}=\boxed{20\,^\circ\mathrm{C}},
      \qquad
      T_2(0)=\frac{30}{2}=\boxed{15\,^\circ\mathrm{C}}.
    \]

    \textit{Criteri:} 0,5 punti per ciascuna temperatura.
  \end{solution}

  \part[2]
  Dopo quanti minuti il primo componente
  raggiunge la temperatura di $40\,^\circ\mathrm{C}$?
  Imposta e risolvi un'equazione.

  \begin{solution}
    Nel contesto $t\ge0$, quindi $t+3\ne0$.
    \[
      \frac{60(t+1)}{t+3}=40
      \quad\Longleftrightarrow\quad
      60t+60=40t+120
      \quad\Longleftrightarrow\quad
      t=3.
    \]

    Il valore è ammissibile: la temperatura richiesta
    si raggiunge dopo $\boxed{3\text{ minuti}}$.

    \textit{Criteri:} equazione e dominio 1 punto;
    risoluzione e risposta 1 punto.
  \end{solution}

  \part[2]
  Dopo quanti minuti i due componenti hanno
  la stessa temperatura? Qual è tale temperatura?
  Imposta e risolvi un'equazione.

  \begin{solution}
    Per $t\ge0$ entrambi i denominatori sono non nulli.
    \begin{align*}
      \frac{60(t+1)}{t+3}
      &=\frac{30(2t+1)}{t+2}\\
      2(t+1)(t+2)&=(2t+1)(t+3)\\
      2t^2+6t+4&=2t^2+7t+3\\
      t&=1.
    \end{align*}

    Il valore è ammissibile e
    \[
      T_1(1)=\frac{60\cdot2}{4}=30,
      \qquad
      T_2(1)=\frac{30\cdot3}{3}=30.
    \]

    Le temperature coincidono dopo $\boxed{1\text{ minuto}}$
    e valgono $\boxed{30\,^\circ\mathrm{C}}$.

    \textit{Criteri:} equazione e passaggi equivalenti 1 punto;
    tempo ammissibile 0,5 punti;
    temperatura 0,5 punti.
  \end{solution}
\end{parts}
\end{esercizio}

\end{esercizi}
\finefila

\end{document}
